\problem{P02}{对称分布的反向配对估计}{对称性与方差}
随机变量 $S$ 与 $-S$ 同分布，记 $\tau=\Prob(S=0)$。为了估计 $p=\Prob(S>0)$，比较两个估计量：用两个独立样本 $S_1,S_2$ 构造
\[
A=\frac{\ind\{S_1>0\}+\ind\{S_2>0\}}2,
\qquad B=\frac{\ind\{S>0\}+\ind\{-S>0\}}2.
\]
证明二者无偏，求其方差，并说明何时反向配对完全消除随机误差。
\solution 对称性给出 $\Prob(S>0)=\Prob(S<0)=(1-\tau)/2$，所以 $\E A=\E B=p$。$A$ 中的两个指示变量独立，故
\[
\Var A=\frac{p(1-p)}2=\frac{1-\tau^2}{8}.
\]
而 $S$ 与 $-S$ 不可能同时为正；除 $S=0$ 外恰有一个为正。因此
\[
B=\frac{1-\ind\{S=0\}}2,
\qquad \Var B=\frac{\tau(1-\tau)}4.
\]
相减得到
\[
\boxed{\Var A-\Var B=\frac{(1-\tau)^2}{8}\ge0.}
\]
当 $0\le\tau<1$ 时，方差之比为 $2\tau/(1+\tau)$。特别地，$\tau=0$ 时 $B\equiv1/2$，配对估计没有随机误差。$\tau=1$ 时两种估计均恒为零；此时不能再以方差之比作比较。
\insight 配对的收益来自负相关。边际分布相同并不意味着两次计算必须独立；这里把正负事件配在一起，只留下零点事件的不确定性。
